\documentclass[10pt,a4paper,twocolumn,twoside]{article}
\usepackage[utf8]{inputenc}
\usepackage[spanish]{babel}
\usepackage{multicol}
\usepackage{graphicx}
\usepackage{fancyhdr}
\usepackage{times}
\usepackage{titlesec}
\usepackage{multirow}
\usepackage{lettrine}
\usepackage[top=2cm, bottom=1.5cm, left=2cm, right=2cm]{geometry}
\usepackage[figurename=Fig.,tablename=TAULA]{caption}
\usepackage{url}

\author{\LARGE\sffamily Luis Martínez Zamora}
\title{\Huge{\sffamily Data-Centric AI\@: Millora de models de diagnòstic mèdic mitjançant pipelines de curació i ingestió automàtica de vídeo}}


\newcommand\blfootnote[1]{%
  \begingroup
  \renewcommand\thefootnote{}\footnote{#1}%
  \addtocounter{footnote}{-1}%
  \endgroup
}

%
%\large\bfseries\sffamily
\titleformat{\section}
{\large\sffamily\scshape\bfseries}
{\textbf{\thesection}}{1em}{}

\begin{document}

\fancyhead[LO]{\scriptsize Luis Martínez Zamora: Data Centric AI\@: Millora de models de diagnòstic mèdic mitjançant pipelines de curació i ingestió automàtica de vídeo}
\fancyhead[RO]{\thepage}
\fancyhead[LE]{\thepage}
\fancyhead[RE]{\scriptsize EE/UAB TFG INFORMÀTICA\@: Data-Centric AI}

\fancyfoot[CO,CE]{}

\fancypagestyle{primerapagina}
{
  \fancyhf{}
  \fancyhead[L]{\scriptsize TFG EN ENGINYERIA INFORMÀTICA, ESCOLA D'ENGINYERIA (EE), UNIVERSITAT AUTÒNOMA DE BARCELONA (UAB)}
  \fancyfoot[C]{\scriptsize ``Mes''~de~2026, Escola d'Enginyeria (UAB)}
}

%\lhead{\thepage}
%\chead{}
%\rhead{\tiny EE/UAB TFG INFORMÀTICA: TÍTOL (ABREUJAT SI ÉS MOLT LLARG)}
%\lhead{ EE/UAB \thepage}
%\lfoot{}
%\cfoot{\tiny{Mes~2026, Escola d'Enginyeria (UAB)}}
%\rfoot{}
\renewcommand{\headrulewidth}{0pt}
\renewcommand{\footrulewidth}{0pt}
\pagestyle{fancy}

%\thispagestyle{myheadings}
\twocolumn[\begin{@twocolumnfalse}

%\vspace*{-1cm}{\scriptsize TFG EN ENGINYERIA INFORMÀTICA, ESCOLA D'ENGINYERIA (EE), UNIVERSITAT AUTÒNOMA DE BARCELONA (UAB)}

\maketitle

\thispagestyle{primerapagina}
%\twocolumn[\begin{@twocolumnfalse}
%\maketitle
%\begin{abstract}
\begin{center}
\parbox{0.915\textwidth}
{\sffamily
\textbf{Resum--}
Resum del projecte, màxim 10 línies.
\\
\\
\textbf{Paraules clau-- } Paraules clau del treball, màxim 2 línies \\
\\
%\end{abstract}
%\bigskip
%\begin{abstract}
\bigskip
\\
\textbf{Abstract--} Versió en anglès del resum 
\\
\\
\textbf{Keywords-- } Versió en anglès de les paraules clau.\\
}


{\vrule depth 0pt height 0.5pt width 4cm\hspace{7.5pt}%
\raisebox{-3.5pt}{\fontfamily{pzd}\fontencoding{U}\fontseries{m}\fontshape{n}\fontsize{11}{12}\selectfont\char70}%
\hspace{7.5pt}\vrule depth 0pt height 0.5pt width 4cm\relax}

\end{center}

%\end{abstract}
\end{@twocolumnfalse}]

\section{Introducción y objetivos}

\blfootnote{$\bullet$ E-mail de contacte: lmzamoraa@gmail.com}
\blfootnote{$\bullet$ Menció: Computació}
\blfootnote{$\bullet$ Treball tutoritzat per: Yael Tudela (Ciències de la Computació)}
\blfootnote{$\bullet$ Curs 2025/26}

\lettrine[lines=3]{L}{a} colonoscopia es un procedimiento médico que tiene como finalidad examinar el interior del intestino grueso y el recto. Este procedimiento se lleva a cabo mediante la introducción por el ano de un instrumento llamado colonoscopio, un tubo flexible equipado con una cámara y una fuente de luz en su extremo. Esta técnica permite la visualización directa y en tiempo real de la mucosa intestinal, convirtiéndose en una herramienta clínica clave para la detección temprana de úlceras, inflamaciones, tumores y, específicamente en el contexto de este trabajo, pólipos. La identificación precisa de estas lesiones precancerosas durante la exploración en vídeo es crítica, ya que su detección temprana previene de forma significativa el desarrollo de cáncer colorrectal.

Actualmente, la evaluación de pólipos durante una colonoscopia se realiza principalmente a través de la observación visual del médico especialista. Sin embargo, esta tarea puede ser desafiante debido a la variabilidad en la apariencia anatómica de las lesiones, las condiciones de iluminación y la presencia de artefactos visuales. En este contexto, el uso de técnicas de inteligencia artificial para asistir en el análisis automatizado de colonoscopias ha ganado un interés significativo. El objetivo principal de este proyecto es el desarrollo de un pipeline automatizado que permita mejorar el resultado respecto a lo conseguido inicialmente gracias al tratamiento realizado sobre los datos.

Para ello, se ha utilizado el enfoque Data-Centric AI\@. A diferencia del enfoque tradicional, que busca mejorar los resultados diseñando arquitecturas de redes neuronales cada vez más complejas, esta aproximación postula que la calidad y representatividad de los datos son el factor limitante en el rendimiento del sistema. Partiendo de un conjunto de datos inicial proporcionado por especialistas médicos, el cual presenta un severo desbalanceo de clases, este proyecto se centra en el diseño y desarrollo de un pipeline de software automatizado para la ingestión de vídeos completos de colonoscopias. Dicho sistema se encarga de la extracción de frames, el filtrado de calidad y la curación de los datos para balancear las clases minoritarias. El éxito de este trabajo radica en demostrar cómo la mejora sistemática de los datos de entrenamiento a través de la ingeniería de software permite que un modelo de clasificación estándar mejore significativamente su rendimiento.

\subsection{Objetivos}
\subsubsection{Objetivo principal}
Desarrollar un pipeline automatizado de curación de datos bajo el paradigma Data-Centric AI, evaluando de forma iterativa la evolución del rendimiento de un modelo estándar de clasificación de pólipos en cada etapa del proceso. El propósito es cuantificar la mejora progresiva del sistema paso a paso y demostrar empíricamente que la optimización de los datos es el factor determinante para el éxito del modelo.
\subsubsection{Objetivos secundarios}
\begin{itemize}
    \item Construir un módulo de software eficiente para el procesamiento de vídeos de colonoscopias que permita la extracción automatizada de frames contiguos a partir de las marcas temporales (timestamps) etiquetadas por los médicos.
    \item Implementar técnicas de Visión Artificial clásica para la evaluación de la calidad de imagen, utilizando métricas como la Varianza del Laplaciano para el desenfoque, análisis en espacios de color (HSV/LAB) y Entropía de Shannon para imágenes muy oscuras o con tonos uniformes. El sistema detectará y descartará automáticamente frames sin valor diagnóstico (imágenes borrosas, oscuras, choques con la mucosa, exceso de lúmen o alta saturación por reflejos especulares), generando métricas de trazabilidad sobre los motivos de exclusión.
    \item Implementar algoritmos de similitud visual, haciendo uso de técnicas como Perceptual Hashing (pHash), Image Matching o métricas como Structural Similarity Index Measure (SSIM), para identificar y descartar de forma automatizada frames altamente redundantes extraídos del vídeo. Esta técnica garantizará variación entre las imágenes intra-clase y prevendrá el sobreajuste (overfitting) del modelo. Posteriormente, sobre este conjunto de datos curado y único, aplicar estrategias de muestreo para equilibrar la representación de las clases minoritarias en el dataset final.
    \item Definir y ejecutar un marco de experimentación que registre y compare las métricas del modelo tras la aplicación sucesiva de cada módulo del pipeline, aislando y demostrando el impacto individual de cada técnica.
\end{itemize}

\section{Estado del arte}
En la actualidad, la integración de la Inteligencia Artificial en el análisis de colonoscopias se clasifica en dos paradigmas principales. Por un lado, los sistemas de Detección Asistida por Ordenador (CADe), cuyo objetivo clínico es la localización espacial de los pólipos en la imagen o durante la exploración en tiempo real. Esta tecnología ya se encuentra desplegada activamente en los quirófanos de diversos centros hospitalarios. Por otro lado, los sistemas de Diagnóstico Asistido por Ordenador (CADx), en los cuales se enmarca este proyecto, abordan un desafío computacional de mayor complejidad semántica: la clasificación y caracterización óptica del pólipo una vez este ha sido detectado.

Dentro del ámbito CADx, gran parte de los proyectos existentes se han limitado a desarrollar clasificadores binarios que discriminan únicamente entre lesiones neoplásicas (adenomas precancerosos) y no neoplásicas (pólipos benignos). Sin embargo, desde un punto de vista clínico, resulta indispensable desarrollar sistemas capaces de clasificar todo el espectro de tipologías de pólipos, ya que cada variante exige un tratamiento distinto. Computacionalmente, esta transición de una clasificación binaria a una multi-clase introduce un problema crítico: el severo desbalanceo natural de las muestras. La escasez de imágenes para las tipologías menos comunes provoca que los clasificadores estándar pierdan eficacia, haciendo indispensable la aplicación de técnicas avanzadas de curación de datos.

\subsection{Clasificación de pólipos}
Mayoritariamente, el desarrollo de sistemas CADx para la clasificación de pólipos se ha basado en el paradigma centrado en el modelo (Model-Centric AI). Bajo este enfoque, la comunidad científica ha tratado de superar los retos de escasez de datos y el desbalanceo de clases mediante el diseño de procesos más complejos computacionalmente.\\
Un claro ejemplo de esta tendencia es el uso de cuatro redes neuronales de forma paralela para clasificar el tamaño de los pólipos\cite{haldar2023xgboosted}, una métrica clínica relevante para evaluar el riesgo de malignidad de la lesión. En un nivel de complejidad similar destaca el sistema PolyDSS, el cual encadena redes pesadas\cite{badawy2024polydss}, un sistema que encadena redes pesadas de segmentación (ResUNet) con modelos de clasificación (EfficientNet), junto a sistemas de votación múltiple (ensemble voting). Aunque ambos enfoques reportan precisiones en laboratorio superiores al 90\%, introducen una barrera técnica crítica para su despliegue clínico: el alto coste computacional que no permite su uso en tiempo real durante la exploración endoscópica.\\
Desde la perspectiva de la Ingeniería de Software médico, la viabilidad de un sistema CADx está estrictamente condicionada a su capacidad para operar en tiempo real (Real-Time). Arquitecturas compuestas como las mencionadas previamente exigen tiempos de inferencia por frame que superan holgadamente esta restricción temporal. Desplegar estos modelos provocaría latencia (lag) en el monitor del quirófano, haciendo inviable su uso rutinario.\\
Por el contrario, otras líneas de investigación han intentado mitigar este problema de latencia optimizando drásticamente la arquitectura del modelo para su despliegue en dispositivos de quirófano como Colon-YOLO\cite{colonyolo2025}. Mediante la modificación intensiva de un modelo de detección de única etapa, el autor logra viabilidad para la práctica clínica.\\
Sin embargo, el análisis de este sistema revela una limitación fundamental del enfoque Model-Centric: a pesar de la extrema optimización arquitectónica, el rendimiento de clasificación sigue estando relacionado por la calidad del dataset. En este estudio, las clases minoritarias ven reducida su precisión debido al desbalanceo severo, un problema que se intenta paliar con técnicas de aumento de datos espacial (Data Augmentation).

\subsection{Data-Centric AI\@}
Para superar la barrera mencionada en el anterior apartado, este trabajo adopta el paradigma Data-Centric AI\cite{zha2025datacentric}. Desde una perspectiva algorítmica, el enfoque tradicional intenta maximizar el rendimiento del sistema iterando sobre la arquitectura del modelo $M$ mientras mantiene estático un conjunto de datos $D$. Por el contrario, la metodología Data-Centric asume una arquitectura $M$ fija, normalmente ligera y eficiente, mientras que traslada toda la carga de ingeniería a optimizar de forma sistemática y automatizada el conjunto de datos $D$. En la práctica, la aplicación de esta metodología requiere abandonar las limpiezas de datos manuales para diseñar pipelines de software estructurados, específicamente orientados a la curación de datos de entrenamiento (Training Data Curation). Este enfoque consiste en transformar datos en bruto en un conjunto de alto valor. Para ello, el sistema debe operar sistemáticamente sobre tres ejes fundamentales: la mitigación del ruido visual, la eliminación de la redundancia temporal y la resolución del desbalanceo de clases.\\
Sin embargo, la mayoría de los estudios médicos que intentan mitigar el desbalanceo o la mala calidad de las imágenes recurren a técnicas de curación manual o se limitan a aplicar técnicas de aumento de datos espacial a nivel de modelo. Se identifica un vacío tecnológico en la integración sistemática de técnicas de Visión Artificial clásica para el filtrado de ruido (oclusiones, reflejos) junto con algoritmos de similitud visual para eliminar la fuga de datos temporal (data leakage) derivada de la extracción de frames. Por consiguiente, el desarrollo de un sistema automatizado que unifique la extracción, la deduplicación y el balanceo de clases se presenta como el paso lógico e indispensable para viabilizar el uso de clasificadores CADx en tiempo real.


\bibliographystyle{IEEEtran}
\bibliography{referencias}

\end{document}
